% GDZ PPN236018647,native506;HKUST507;original diary9r,service19.
% Complete manuscript entries corresponding to printed108-113.
\clearpage\SourceLeaf{506}\NoSourceHead
\begingroup\fontsize{10}{15}\selectfont\itshape
\newcommand{\DiaryCancelSix}[1]{\tikz[baseline=(t.base)]{\node[inner sep=0pt,outer sep=0pt] (t) {\strut #1};\draw[line width=.7pt] ([xshift=-.06em,yshift=.05ex]t.west)..controls([yshift=.25ex]t.center)..([xshift=.08em,yshift=-.08ex]t.east);}}
\noindent\tikz[baseline=-.25ex,x=.85em,y=.85em]{
\draw[line width=.5pt] (.15,-.65)--(.62,1.05);\draw[line width=.4pt] (.34,-.68)--(.86,1.04);\draw[line width=.4pt] (.55,-.7)--(1.03,1.03);\draw[line width=.35pt] (.78,-.65)--(1.25,1.02);\draw[line width=.4pt] (.98,-.65)--(1.36,.82);
\draw[line width=.6pt] (.19,.7)--(1.48,.56);\draw[line width=.5pt] (.13,.42)--(1.34,.33);\draw[line width=.45pt] (.13,.18)--(1.23,.04);\draw[line width=.5pt] (.09,.04)--(1.35,-.14);\draw[line width=.45pt] (.12,-.18)--(1.34,-.28);
\draw[line width=.4pt] (1.62,-.25)--(2.08,1.01);\draw[line width=.35pt] (1.66,-.64)--(2.25,1.2);\draw[line width=.35pt] (1.94,-.69)--(2.4,1.08);\draw[line width=.35pt] (2.05,-.59)--(2.54,1.16);\draw[line width=.4pt] (2.2,-.34)--(2.65,1.2);\draw[line width=.35pt] (2.28,-.23)--(2.74,1.07);
\draw[line width=.45pt] (1.8,.84)--(2.96,.46);\draw[line width=.45pt] (1.82,.5)--(2.94,.25);\draw[line width=.5pt] (1.54,.41)--(2.88,.09);\draw[line width=.6pt] (1.73,.16)--(2.89,-.09);\draw[line width=.45pt] (1.83,.07)--(2.56,-.26);
}\quad Numeratorem et denominatorem sinus lemni-\\
scatici (universalissime accepti) ad quantitates integra-\\
les reducere contigit; simul omnium functionum\\
lemniscaticarum quae excogitari possunt, evolutiones\\
in series infinitas ex principiis genuinis haustae,\\
inventum pulcherrimum sane nullique praecedentium inferius

\medskip
Praeterea iisdem diebus principia deteximus secundum\\
quae series arithmetico-geometricae interpolari debent, ita\\
ut \DiaryCancelSix{omnes} terminos in progressione data ad \DiaryCancelSix{term} indicem quem-\\
cumque rationalem pertinentes per aequationes algebraicas\\
exhibere iam in potestate sit.\hfill Mai. ult. Jun 2. 3.

\medskip
\noindent\tikz[baseline=-.3ex,x=.8em,y=.8em]{
\draw[line width=.35pt] (.02,-.65)--(.61,.65);\draw[line width=.35pt] (.12,-.65)--(.76,.69);\draw[line width=.35pt] (.32,-.6)--(.98,.67);
\draw[line width=.5pt] (.31,.65)--(1.01,-.02);\draw[line width=.45pt] (.19,.35)--(.86,-.3);\draw[line width=.45pt] (.07,.05)--(.72,-.51);
\draw[line width=.35pt] (1.22,-.55)--(1.78,.6);\draw[line width=.35pt] (1.37,-.62)--(1.94,.67);\draw[line width=.35pt] (1.57,-.58)--(2.16,.65);
\draw[line width=.5pt] (1.49,.67)--(2.06,.02);\draw[line width=.45pt] (1.38,.41)--(1.97,-.27);\draw[line width=.45pt] (1.29,.09)--(1.84,-.52);
}\quad Inter duos numeros datos \tikz[baseline=(lo.base)]{
\node[inner sep=0pt,outer sep=0pt] (lo) {\strut duo};
\node[inner sep=0pt,outer sep=0pt,anchor=south] (hi) at ([xshift=-.4em,yshift=.4ex]lo.north) {\strut ut};
\draw[line width=.5pt] ([xshift=-.03em,yshift=.05ex]lo.west)--([xshift=.12em,yshift=-.05ex]lo.east);
\draw[line width=.3pt] ([yshift=-.22ex]lo.south west)--([yshift=-.1ex]lo.south east);
\draw[line width=.5pt] ([xshift=-.04em,yshift=.1ex]hi.west)..controls([yshift=.3ex]hi.center)..([xshift=.22em,yshift=.06ex]hi.east);
} semper dantur\\
infinite multi termini medii tum arithmetico geometrici tum\\
harmonico geometrici, quorum nexum mutuum ex asse\\
perspiciendi felicitas nobis est facta.\hfill Junio 3. Brunov.

\medskip
\noindent\tikz[baseline=-.3ex,x=.7em,y=.7em]{
\draw[line width=.3pt] (.02,-.4)--(.7,.7);\draw[line width=.3pt] (.26,-.57)--(.99,.85);
\draw[line width=.45pt] (.27,.67)..controls(.4,.23)and(.62,.12)..(.83,.04);
\draw[line width=.45pt] (.04,.16)..controls(.35,-.1)and(.5,-.12)..(.75,-.15);
\draw[line width=.4pt] (0,-.03)..controls(.25,-.47)and(.48,-.41)..(.7,-.46);
}\quad Theoriam nostram iam ad Transcendentes\\
ellipticas immediate applicavimus\hfill Junio 5

\medskip
\noindent\tikz[baseline=-.3ex,x=.6em,y=.6em]{
\draw[line width=.35pt] (.12,-.5)--(.62,.62);\draw[line width=.3pt] (.3,-.54)--(.86,.64);
\draw[line width=.4pt] (.37,.55)--(1.04,.35);\draw[line width=.4pt] (.23,.27)--(.89,.11);\draw[line width=.4pt] (.11,-.04)--(.75,-.13);
}\quad Rectificatio Ellipseos tribus Modis di-\hfill Juni 10\\
versis absoluta.

\medskip
\noindent\tikz[baseline=-.3ex,x=.6em,y=.6em]{
\draw[line width=.5pt] (.24,-.5)--(.3,.79);\draw[line width=.4pt] (.42,-.57)--(.54,.74);\draw[line width=.4pt] (.59,-.7)--(.78,.62);
\draw[line width=.5pt] (0,.45)--(.81,.32);\draw[line width=.5pt] (0,.14)--(.77,.01);\draw[line width=.5pt] (.02,-.18)--(.68,-.22);
}\quad Calculum Numerico-Exponentialem omnino no-\\
vum invenimus\quad\tikz[baseline=-.3ex,x=1em,y=1ex]{
\draw[line width=.35pt] (0,.15)--(.55,.1);\draw[line width=.6pt] (1.2,.2)--(1.8,.13);\draw[line width=.4pt] (2.5,.22)--(2.82,.14);\draw[line width=.5pt] (3.6,.45)--(3.82,.05);\draw[line width=.45pt] (4.4,.4)--(4.72,.28);
}\hfill Jun 12

\medskip
\noindent Problema e calculo probabilitatis circa fractiones continuas olim\\
frustra tentatum solvimus\hfill Oct. 25
\endgroup
